\phantomsection\label{sec:construction-guide}

Part~\ref{part:main} states what the limiting theory must satisfy.
Here we explain how to construct it from the Wilson measure of Definition~\ref{def:wilson-law}.
The construction first removes the spacing at each fixed volume, then takes a subsequence as the volume grows. The later sections explain how those limiting correlations define fields, symmetries and phases.

\paragraph{Following an observable across scales.}
Section~\ref{sec:uv-lattice} starts with the finite probability law and its exact symmetries.
It describes a classical local observable by a gauge-invariant polynomial in the Taylor coefficients of a smooth connection (Definition~\ref{def:uv-labels}); the connection specifies infinitesimal parallel transport, and these coefficients record its behaviour near one point.
A polynomial in nearby link matrices represents that classical polynomial on the lattice.
Integrating out fine links can change its coefficients, so Definition~\ref{def:observables} gives it a fixed index and~\eqref{eq:uv-insertion} fixes the coefficients at a reference level.
This lets us follow one observable through different lattice spacings.

Section~\ref{sec:uv-activities} constructs the link polynomials and organises products of insertions for integration.
It records each contribution together with the lattice region on which it depends; terms where a plaquette is far from the identity are kept and tracked by their sites.
This bookkeeping prepares the estimates but does not prove them.

The later ultraviolet sections perform these integrations one spacing at a time.
They average over finer link variables to produce a density for the remaining coarse links, preserving the original integral.
The estimates must control both this density and each insertion after many such steps.
When every plaquette matrix in a region is close to the identity, the Wilson energy is compared with its quadratic approximation.
When a plaquette is farther away, its contribution is bounded using the energy instead of being discarded.
Contributions are grouped by the sites on which they depend, which is what allows bounds for one local observable to be independent of the surrounding box.
The change in the coefficient of the quadratic curvature term is included in the coupling assigned to the original lattice measure.

These arguments must provide two kinds of estimates.
At a fixed box size, they must show that correlations approach a limit along the full sequence of spacings, as in Proposition~\ref{prop:uv}.
Uniform bounds on local observables and on their finite differences as an insertion point moves must also hold independently of box size.
The latter bounds give a convergent subsequence as the box grows; the fixed-volume limits identify its inner limits with those of the full spacing sequence.
For the moments of Definition~\ref{def:moments}, the intended passage is the ordered limit~\eqref{eq:ordered}:
\begin{equation*}
\begin{aligned}
 C_{N,p}(v;f)&\longrightarrow R_p &&(N\to\infty,\ p\text{ fixed}),\\
 R_{p(q(k))}&\longrightarrow\int f(x)S^\mix(\pi(v_1),\ldots,\pi(v_n);x)\,dx
 &&(k\to\infty).
\end{aligned}
\end{equation*}
Here $f$ is a separated test of Definition~\ref{def:tests}; the inner limit is required at every volume containing its support (Definition~\ref{def:admissible}).
The family is chosen before the prescribed volume sequence $p$, and the subsequence $q$ is the same for every observable list and test.

\paragraph{From correlation functions to fields.}
Section~\ref{sec:infinite-volume} applies reflection positivity: for positive-time products $F,G$, reflect $F$ across time zero and take its complex conjugate, writing the result as $\vartheta F$.
The resulting positive form is $S^{\mix}((\vartheta F)G)\ge0$; the finite-lattice inequality is~\eqref{eq:iv-lattice-rp}, and its repeated-cell version is~\eqref{eq:iv-chessboard}.
Repeated reflection bounds a joint correlation by one-cell bounds, which do not grow with the box.
These bounds give one subsequence for all correlations.
To identify derivatives, the proof compares a finite difference of an insertion's position with a higher-degree observable.
Two observables are identified as the same continuum field when every tested limiting correlation of their difference is zero; removing this null subspace gives the field spaces of the limit.
Smoothness therefore needs estimates for these comparisons, not just compactness.
Finally the fixed-volume theorem removes the temporary spacing subsequence and gives the ordered limit of Definition~\ref{def:ordered}.

Section~\ref{sec:maxwell} compares the limiting fields at short distances with the explicit curvature fields of the free Maxwell theory.
These are the field strengths of a free electromagnetic field; finite tests distinguish the field classes, and comparison maps transfer the stress-tensor estimates to the interacting theory.
Constructing these maps requires short-distance bounds and a separate dimension bound.
The comparison does not make the interacting theory free.
For the symmetry argument, the stress-tensor integral identities are a separate required input to establish continuous rotations on this space; they are not proved in the sections presented here.

\paragraph{Reading the construction.}
When reading an estimate, identify the quantities held fixed, the limit being taken, and whether its bound depends on volume.
This is especially important when a fixed-volume spacing limit is used inside the later volume subsequence.

